\section{Introduction}
\label{sec:intro}
A language model does one thing: given a sequence of tokens, it produces a distribution over the next token. It cannot see a file it was not shown, remember a conversation it was not sent again, or carry out a command it wrote. Language models are nonetheless described as agents that read codebases, run tests and fix bugs. The difference between what a model is and what an agent does is software around the model, and that software is the subject of this paper.

The software has many names. Evaluations call it a \emph{scaffold} \cite{anthropic2025swebench,metr2025update}; libraries call themselves \emph{frameworks}, \emph{SDKs} or an \emph{orchestration framework and runtime} \cite{langgraph}; and \emph{harness} is now common among people who build coding agents \cite{anthropic2025harnesses}. None of the names says what the thing contains, where it ends and the model begins, or how two instances of it differ. Agents are therefore compared by the model inside them, although results vary with the system around it (\S\ref{sec:uses}), and when an agent fails there is no agreed list of places to look.

\paragraph{Position} We argue three claims.
\begin{enumerate}
  \item \textbf{An LLM agent is a model wrapped in a harness}: $\textit{Agent} = \textit{Harness}(\textit{Model})$. The model maps tokens to a next token. The harness is everything else on the path of a request.
  \item \textbf{Every harness is built from five mechanistic primitives}: \prim{Input}, \prim{Context}, \prim{Model interface}, \prim{Output} and \prim{Control flow}. Each is defined by what it does and by what it never does, so that each piece of a harness has one home.
  \item \textbf{Everything else said about a harness is at one of two other levels}: a \emph{design decision} within a primitive (one tool or forty; a loop or a graph), or a \emph{production concern} asked of every primitive (persistence, safety and security, observability, resilience, performance, evaluation). Neither adds a primitive.
\end{enumerate}
The vocabulary separates what a harness \emph{is} from what a given harness \emph{chose} and from what makes it fit to \emph{run}, so that a failure can be located in one of five places before it is attributed to the model, and a benchmark result can report the harness that produced it.

\paragraph{Evidence} The definitions come from building a harness from nothing, one primitive at a time, around a fixed model and task, and observing what each version could not do; the five programs are published with the paper. The claim that the five cover existing harnesses is tested by deconstructing six of them from their primary sources at pinned versions (\S\ref{sec:deconstruction}): every component on the request path is assigned to one primitive, or to one concern together with the primitive it changes, and every case that was hard to place is recorded. The evidence, a cited assignment for each of 548 components, is published with the paper. The full paper, with the per-harness deconstructions, the rules in full and the appendices, is available as a preprint; this paper is its summary, shared non-archivally. This is a conceptual contribution: it reports no experiment in which a primitive is varied, and it does not prove that no sixth primitive could ever be needed. It states what would refute the claim (\S\ref{sec:limitations}) and gives a deconstruction a reader can check.
